\documentclass[11pt,a4paper]{article}

\usepackage[margin=1in]{geometry}
\usepackage{amsmath}
\usepackage{amssymb}
\usepackage{booktabs}
\usepackage{longtable}
\usepackage{graphicx}
\usepackage{caption}
\usepackage[hidelinks]{hyperref}

\title{Monte Carlo Scenario Generation for the SBA Asset--Liability Model\\
       \large Diagnostic Findings and Algorithm Selection}
\author{Prepared for the SBA project}
\date{27 September 2026}

\newcommand{\bp}{\,\text{bp}}

\begin{document}
\maketitle

\begin{abstract}
This note selects a Monte Carlo algorithm for generating forward scenarios from
the SBA term-structure model. Rather than choosing on theoretical grounds, we
estimate the existing model on the full sample (284 monthly observations,
January 2003 -- August 2026) and let its residual diagnostics drive the choice.
Three findings determine the recommendation. First, the dominant credit-spread
factor has excess kurtosis near $10$ and positive skewness, which disqualifies
Gaussian innovations. Second, information criteria disagree sharply on lag
order---AIC selects eight lags against BIC's one---and the parsimonious model is
the appropriate one for simulation. Third, and most consequentially, applying
first differences to the credit-spread factors discards their mean reversion and
produces scenario sets in which $47\%$ of paths contain an economically
impossible negative credit spread, against a historical base rate of $3.5\%$. We
recommend a semiparametric VAR-bootstrap on a mixed-integration-order state
vector, with the spread curve reparameterised as a shifted logarithm. This
reduces the violation rate to $4.2\%$, consistent with history.
\end{abstract}

\section{Context and objective}

The SBA project builds an economic scenario generator whose terminal purpose is
portfolio allocation: scenarios feed estimates of expected return, risk,
correlation and macroeconomic state, which in turn drive allocation decisions
against a mortality-driven liability. The scenario set is therefore an
intermediate product, and its quality is judged by whether allocation decisions
taken against it are sound---not by in-sample fit alone.

Two requirements follow, and they are not the same requirement:

\begin{description}
  \item[Statistical plausibility.] Simulated paths should reproduce the
    distributional and dependence properties of the historical data: fat tails,
    skewness, cross-factor correlation, and volatility clustering.
  \item[Economic coherence.] Simulated states must be states the economy can
    actually occupy. A scenario set can be well calibrated in every moment and
    still be unusable if it contains negative credit spreads.
\end{description}

Section~\ref{sec:coherence} shows that the current model satisfies the first
requirement tolerably and fails the second badly.

\subsection{Model as estimated}

The pipeline reduces each of two curves to three principal components, stacks the
six scores, and fits a vector autoregression to their first differences:

\begin{equation}
  \Delta \mathbf{f}_t = \mathbf{c} + \sum_{i=1}^{p} \mathbf{A}_i \,
  \Delta \mathbf{f}_{t-i} + \boldsymbol{\varepsilon}_t ,
  \qquad
  \mathbf{f}_t = \big(\text{TPC}_{1:3,t},\ \text{SPC}_{1:3,t}\big)^{\!\top} ,
\end{equation}

where the Treasury scores $\text{TPC}$ derive from the Treasury nominal spot
curve and the spread scores $\text{SPC}$ from the HQM corporate spot curve minus
the Treasury curve, both on a common 200-point maturity grid spanning
$0.5$ to $100$ years.

\section{Factor structure}

Three components suffice for both curves. The Treasury curve retains $99.83\%$
of variance and the spread curve $99.45\%$, with the conventional level, slope
and curvature interpretation.

\begin{table}[htbp]
\centering
\caption{Explained variance ratio by component.}
\label{tab:evr}
\begin{tabular}{lcccc}
\toprule
Curve & PC1 & PC2 & PC3 & Cumulative \\
\midrule
Treasury      & 0.9063 & 0.0821 & 0.0099 & 0.99829 \\
Credit spread & 0.8703 & 0.0980 & 0.0262 & 0.99445 \\
\bottomrule
\end{tabular}
\end{table}

The discarded $0.17\%$ is not uniformly distributed across maturities, and this
matters more than the headline figure suggests. Reconstructing the Treasury
curve from three components gives a root-mean-square error of $4.5\bp$ overall
but $26.7\bp$ at the six-month point, with a maximum absolute error of
$99.8\bp$.

\begin{table}[htbp]
\centering
\caption{Three-component reconstruction error for the Treasury curve, by maturity.}
\label{tab:recon}
\begin{tabular}{lrrrrrrr}
\toprule
Maturity & 0.5y & 1y & 2y & 5y & 10y & 30y & 100y \\
\midrule
RMSE (bp) & 26.70 & 16.65 & 6.32 & 12.71 & 4.58 & 4.12 & 1.60 \\
\bottomrule
\end{tabular}
\end{table}

Any simulation that reconstructs curves from simulated factors must therefore
decide explicitly what to do with the residual. Discarding it silently
introduces a systematic error of roughly a quarter of a percentage point at the
short end, which propagates directly into short-horizon bond returns and into
the discounting of near-dated liability cash flows.

\section{Residual diagnostics}
\label{sec:resid}

\subsection{Lag order}

Information criteria disagree materially. AIC selects eight lags; BIC and HQIC
both select one. The difference is not cosmetic: the AIC specification estimates
$49$ parameters per equation ($294$ in total) against $275$ usable observations,
whereas the BIC specification estimates seven.

\begin{table}[htbp]
\centering
\caption{Lag selection and companion-matrix stability. The largest eigenvalue
modulus governs whether simulated paths remain bounded over long horizons.}
\label{tab:lags}
\begin{tabular}{lrrrr}
\toprule
Criterion & Lags & Parameters/eq. & Observations & $\max|\lambda|$ \\
\midrule
AIC        & 8 & 49 & 275 & 0.9089 \\
BIC        & 1 &  7 & 282 & 0.4318 \\
HQIC       & 1 &  7 & 282 & 0.4318 \\
\bottomrule
\end{tabular}
\end{table}

Both specifications are stable, so neither will produce explosive paths. The
choice rests on estimation variance instead. Over a sixty-month simulation
horizon, coefficient error compounds, and the heavily parameterised AIC model
carries far more of it for no diagnostic benefit: the portmanteau whiteness test
rejects at $p < 0.001$ under \emph{both} specifications. Adding seven lags does
not purchase white residuals, and so we prefer the parsimonious model and handle
the residual dependence in the resampling scheme instead.

\subsection{Distributional properties}

The residuals are decisively non-Gaussian in two of six factors, and the two are
economically important ones.

\begin{table}[htbp]
\centering
\caption{Residual diagnostics for the BIC-selected specification. JB denotes the
Jarque--Bera normality test; ARCH denotes Engle's test for conditional
heteroskedasticity at twelve lags. Kurtosis is reported raw, so $3$ is the
Gaussian reference value.}
\label{tab:diag}
\begin{tabular}{lrrrr}
\toprule
Factor & JB $p$ & ARCH $p$ & Skewness & Kurtosis \\
\midrule
TPC1 & $1.1\times10^{-1}$   & $9.4\times10^{-1}$   & $-0.02$ & 3.62 \\
TPC2 & $6.2\times10^{-1}$   & $2.7\times10^{-2}$   & $-0.10$ & 3.21 \\
TPC3 & $1.1\times10^{-44}$  & $2.0\times10^{-1}$   & $-0.80$ & 6.89 \\
SPC1 & $7.5\times10^{-136}$ & $9.1\times10^{-2}$   & $ 1.17$ & 9.99 \\
SPC2 & $2.4\times10^{-1}$   & $3.6\times10^{-3}$   & $-0.04$ & 3.49 \\
SPC3 & $6.8\times10^{-1}$   & $1.7\times10^{-1}$   & $-0.00$ & 3.26 \\
\bottomrule
\end{tabular}
\end{table}

SPC1, the dominant credit-spread factor, has kurtosis close to $10$ and
\emph{positive} skewness of $1.17$. This is the correct economic signature:
credit spreads widen abruptly and tighten slowly. Gaussian innovations cannot
represent an asymmetry of this kind, and a scenario set built on them will
understate precisely the credit stress episodes that dominate the loss
distribution of a bond-heavy portfolio.

Conditional heteroskedasticity is present but not pervasive, reaching
significance at the $5\%$ level in TPC2 and SPC2 only. This is relevant to the
choice between a simple residual bootstrap and a GARCH-filtered scheme, and we
return to it in Section~\ref{sec:algo}.

\section{Over-differencing and the coherence failure}
\label{sec:coherence}

\subsection{Diagnosis}

Augmented Dickey--Fuller tests on the individual factor levels show that the
six factors are not alike in their persistence.

\begin{table}[htbp]
\centering
\caption{Stationarity of factor levels and implied mean-reversion half-lives.
ADF tests the null of a unit root; KPSS tests the null of stationarity.}
\label{tab:adf}
\begin{tabular}{lrrr}
\toprule
Factor & ADF $p$ (level) & KPSS $p$ (level) & Half-life (months) \\
\midrule
TPC1 & 0.3250 & 0.0100 & 43.0 \\
TPC2 & 0.0932 & 0.0238 & 47.6 \\
TPC3 & 0.0346 & 0.0100 & 13.9 \\
SPC1 & 0.0530 & 0.0100 & 13.1 \\
SPC2 & 0.1115 & 0.0100 & 12.7 \\
SPC3 & 0.2363 & 0.0547 & 15.8 \\
\bottomrule
\end{tabular}
\end{table}

The Treasury level factor TPC1 behaves like an integrated process, with an ADF
$p$-value of $0.33$ and a half-life of $43$ months; differencing it is correct.
The spread factors mean-revert roughly three times faster, with half-lives near
thirteen months, and the ADF test is borderline for SPC1 at $p = 0.053$.
Differencing them imposes a random walk on a series that demonstrably returns to
a central level.

The consequence is visible directly in the spread itself. At the five-year
point, the spread currently stands at $0.56\%$ against a sample mean of $0.99\%$
and a sample minimum of $0.34\%$. Its monthly difference has standard deviation
$0.205\%$, so a random walk implies a sixty-month standard deviation of
$1.588\%$---roughly three times the level of the series being diffused. Nothing
prevents such a process from crossing zero, and in simulation it does so
constantly.

\subsection{The historical base rate is not zero}

One nuance shapes the remedy. The HQM corporate curve does occasionally sit
\emph{below} the Treasury curve in the historical record: $11$ of $56{,}800$
observed cells are non-positive, a rate of $0.019\%$, with a minimum of
$-0.16\%$. Every one of these occurs at the six-month maturity, in eleven months
clustered in 2020 and 2022---episodes of money-market dislocation in which the
short end of a high-quality corporate curve can invert against Treasuries.

This is a genuine feature of the data, not an error. The remedy must therefore
not impose a hard positivity floor. The target is the historical rate of roughly
$3.5\%$ of months containing a negative spread somewhere on the curve, with
essentially none away from the six-month point.

\subsection{Magnitude of the failure}

Table~\ref{tab:coherence} compares the violation rate under the current
specification against two candidate remedies, all evaluated at the sixty-month
horizon with $10{,}000$ paths, a BIC-selected VAR, and identical bootstrap
shocks under a fixed seed.

\begin{table}[htbp]
\centering
\caption{Probability that a simulated path contains a negative credit spread at
month 60. Model A retains the spread factors in levels; Model B additionally
reparameterises the spread curve as $\log(s + \delta)$ with $\delta = 0.5$.}
\label{tab:coherence}
\begin{tabular}{lrr}
\toprule
Specification & $P(\text{spread} < 0)$ & Excluding 0.5y \\
\midrule
Historical base rate                  & 3.52\%  & 0.00\% \\
\midrule
Current: all factors differenced      & 47.08\% & --- \\
Model A: spread factors in levels     & 18.89\% & 14.86\% \\
Model B: shifted-log spread, levels   & 4.21\%  & 2.05\% \\
\bottomrule
\end{tabular}
\end{table}

Under the current specification almost half of all scenarios are economically
impossible. Retaining the spread factors in levels recovers their mean reversion
and cuts the rate to $19\%$; adding the shifted-logarithm reparameterisation
brings it to $4.2\%$, close to the historical $3.5\%$ and concentrated at the
short end where history places it.

\subsection{Tail behaviour under Gaussian and bootstrap innovations}

The distributional finding of Section~\ref{sec:resid} has a measurable cost.
Holding the fitted model fixed and varying only the innovation scheme, the
five-year spread distribution at month 60 differs materially in the tail.

\begin{table}[htbp]
\centering
\caption{Five-year credit spread at month 60 (percentage points), Gaussian
innovations versus stationary block bootstrap of residual vectors.}
\label{tab:tails}
\begin{tabular}{lrrrr}
\toprule
Innovations & Mean & $p_{99}$ & $p_{99.9}$ & Maximum \\
\midrule
Gaussian                    & 0.55 & 5.55 & 7.17 & 8.37 \\
Stationary block bootstrap  & 0.59 & 6.06 & 8.39 & 10.81 \\
\bottomrule
\end{tabular}
\end{table}

The bootstrap produces a far tail roughly $17\%$ wider at the $99.9$th
percentile. For a risk measure defined on the tail---conditional value at risk,
say---this difference flows directly into the allocation.

\section{Recommended algorithm}
\label{sec:algo}

We recommend a \emph{semiparametric VAR-bootstrap}: a parametric VAR capturing
linear dynamics, combined with nonparametric resampling of its residuals. This
is the approach that comparative work on curve scenario generation finds
superior to purely nonparametric resampling where the dynamic model is
adequately specified~\cite{scenariocomp}. Concretely:

\begin{enumerate}
  \item \textbf{Reparameterise the spread curve} as $\log(s_\tau + \delta)$ with
    $\delta = 0.5$. The shift accommodates the genuine historical negatives at
    the short end while the logarithm supplies the asymmetry that makes large
    widenings more likely than large tightenings.

  \item \textbf{Adopt a mixed integration order.} Difference the Treasury
    factors, which are empirically $I(1)$; retain the transformed spread factors
    in levels so that mean reversion is preserved. The state vector becomes
    $\big(\Delta\text{TPC}_{1:3},\ \text{LPC}_{1:3}\big)$.

  \item \textbf{Select lag order by BIC}, not AIC, giving a VAR(1). The AIC
    specification buys no improvement in residual whiteness at more than seven
    times the parameter count.

  \item \textbf{Resample residual vectors by stationary
    bootstrap}~\cite{politisromano} with geometric block lengths of mean six
    months. Resampling whole rows preserves cross-factor dependence; resampling
    in blocks preserves both the serial dependence that the whiteness test
    detects and the volatility clustering that Engle's test finds in TPC2 and
    SPC2, without committing to a parametric volatility model.

  \item \textbf{Restore the PCA residual} when reconstructing curves, either by
    carrying the last observed residual forward or by resampling it alongside
    the factor shocks. Table~\ref{tab:recon} gives the cost of omitting it.

  \item \textbf{Incorporate parameter uncertainty} once the above is in place,
    by bias-corrected bootstrap-after-bootstrap~\cite{kilian}. At a sixty-month
    horizon the sampling error in $\mathbf{A}$ is not negligible, and ignoring
    it makes the predictive distribution too narrow---an error that biases
    allocation toward excessive risk.
\end{enumerate}

\subsection{Alternatives considered and rejected}

\begin{description}
  \item[Gaussian innovations.] Rejected on the evidence of
    Table~\ref{tab:diag}: kurtosis near $10$ and skewness of $1.17$ in the
    dominant spread factor. Retained in the implementation only as a diagnostic
    baseline.

  \item[Filtered historical simulation.] Fitting GARCH models to each factor,
    bootstrapping standardised residuals and reinflating by simulated volatility
    is the industry-standard treatment for volatility
    clustering~\cite{fhs}. Here it is not yet warranted: ARCH effects reach
    significance in only two of six factors, and the block bootstrap addresses
    them nonparametrically. This is the natural upgrade if the scenario set
    later proves to understate clustered stress.

  \item[Quasi-Monte Carlo.] Low-discrepancy sequences accelerate convergence for
    smooth functionals but interact poorly with resampling-based innovations and
    offer least benefit in the deep tail, which is where allocation decisions
    are most sensitive.

  \item[Markov-switching VAR.] Attractive in principle, since the sample spans
    both a zero-interest-rate regime and a normal one. With $284$ monthly
    observations the regime parameters are unlikely to be identified reliably.
\end{description}

\subsection{Implementation notes}

Two practical constraints deserve recording.

\emph{Memory.} Factor paths for $10{,}000$ scenarios over $60$ months occupy
approximately $29$ MB. Materialising the corresponding full $200$-point curves
would require about $1.0$ GB. Simulation should therefore retain factor paths and
reconstruct curves only at the horizons actually required.

\emph{Reproducibility.} All randomness should flow from a single explicitly
threaded generator seeded once, rather than from global state, so that scenario
sets are reproducible and so that competing allocation rules can be compared
under common random numbers.

\section{Connection to the allocation objective}

The recommended scheme produces an equally weighted ensemble of paths. This
composes directly with scenario-based conditional value-at-risk optimisation,
which Rockafellar and Uryasev~\cite{rockafellar} showed can be expressed as a
linear program over exactly such an ensemble, with the value at risk determined
endogenously rather than supplied in advance. The pairing is convenient: no
covariance matrix need be estimated, the fat tails established in
Section~\ref{sec:resid} are carried through to the risk measure intact, and
allocation constraints enter as linear side conditions.

Should narrative macroeconomic scenarios later be assigned subjective
probabilities, those weights must be carried explicitly and must not be mixed
with the equally weighted Monte Carlo ensemble within a single statistic.

\section{Limitations}

Several caveats bound the above.

\begin{itemize}
  \item The diagnostics rest on a single realised path of history spanning
    $2003$--$2026$. The sample contains two severe credit episodes, which
    dominate the estimated tail of SPC1; a different sample window would give
    materially different kurtosis.

  \item ADF and KPSS disagree on the spread factors---ADF marginally rejects a
    unit root for SPC1 while KPSS rejects stationarity. This is the
    characteristic signature of a near-unit-root process, and the levels
    specification is justified on economic grounds as much as statistical ones.

  \item The shift parameter $\delta = 0.5$ is chosen to accommodate the observed
    historical minimum with margin. It has not been estimated, and results
    should be checked for sensitivity to it.

  \item Negative simulated Treasury rates remain at roughly $7\%$ of paths under
    every specification tested. United States nominal rates did not go negative
    in this sample, so this is the next coherence gap to close, most naturally
    by a shadow-rate floor.

  \item No out-of-sample validation of the scenario set has been performed. A
    rolling-origin backtest, in which the model is re-estimated using only data
    available at each decision date, is a prerequisite before any allocation
    conclusion is drawn---particularly because the principal components are
    currently fitted on the full sample, which would constitute look-ahead in a
    backtest.
\end{itemize}

\begin{thebibliography}{9}

\bibitem{fhs}
Barone-Adesi, G., Giannopoulos, K., and Vosper, L. (1999).
\newblock VaR without correlations for portfolios of derivative securities.
\newblock \emph{Journal of Futures Markets}, 19(5), 583--602.

\bibitem{kilian}
Kilian, L. (1998).
\newblock Small-sample confidence intervals for impulse response functions.
\newblock \emph{The Review of Economics and Statistics}, 80(2), 218--230.

\bibitem{littermanscheinkman}
Litterman, R. and Scheinkman, J. (1991).
\newblock Common factors affecting bond returns.
\newblock \emph{Journal of Fixed Income}, 1(1), 54--61.

\bibitem{politisromano}
Politis, D. N. and Romano, J. P. (1994).
\newblock The stationary bootstrap.
\newblock \emph{Journal of the American Statistical Association}, 89(428),
  1303--1313.

\bibitem{rockafellar}
Rockafellar, R. T. and Uryasev, S. (2000).
\newblock Optimization of conditional value-at-risk.
\newblock \emph{Journal of Risk}, 2(3), 21--41.

\bibitem{scenariocomp}
Scenario generation for time series and curves: a comparison of nonparametric
  and semiparametric bootstrap.
\newblock arXiv preprint \texttt{arXiv:2606.11859}, 2026.
\newblock \url{https://arxiv.org/abs/2606.11859}

\end{thebibliography}

\end{document}
